as shown before, despite it's potential the Final Degree model doesn't allow for a better competitiveness. However, unlike in the Semi-Online model, the class of graphs that can force a pessimistic matching of size $\frac{n}{2}$ is very limited. For this reason, it is alwo worth analyzing performance in a much more forgiving, and realistic, enviorenment.

\begin{definition}
	\label{CLV-B Model}
	The CLV-B model is a bipartite adaptation of the CLV model for general graphs introduced by Fan Chung, Linyuan Lu, and Van Vu. It is defined by two vectors, $p\in[0,1]^n,q\in[0,1]^m$.An edge between $u_j\in U,v_i\in V$ appears with probability $p_jq_i$. 
	We assume WLOG that elements of $p$ are in non-increasing order.\\
	We denote the random graph of vectors $p,q$ $I_{p,q}$
\end{definition}
Do note that in this model, much like in the i.i.d model, the adversary is severely resticted and this model resembles a performance on a random graph.
Adapting the RED algorithm to the CLV-B model, with consideration to the fact that there is only the expected value of any final degree and no definitive final degree, results in the following algorithm, desctibed in \cite{obmDegree}
\begin{algorithm}
	\caption{MinPredictedDegree - MPD}
	\label{alg:effective-deg-rand-clv}
	\begin{algorithmic}
		\item on Initialization: for each offline vertex $u_i$ let $d_{u_i}=\sum\limits_{v=\in\vec{v_i}}v\vec{u}_i$
		\item on Arrival of vertex $v$: match to $\underset{u\in U,\text{unmatched}}{\mathrm{argmin}}d_u$
	\end{algorithmic}
\end{algorithm}
\begin{theorem}
	\label{thm:deg-opt}
	MPD is optimal in CLV-B model, meaning\\
	for any algorithm $\mathrm{A}, t\in\mathbb{Z}_{+},p\in[0,1]^n,q\in[0,1]^m$ $\Pr[\mathrm{A}(I_{p,q})>t]\leq \Pr[\MPD(I_{p,q})>t]$
\end{theorem}

the proof, as well as the algorithm itself if from \cite{obmDegree} whose work concerned a slightly different model, where the algorithm has (possibly imperfect) degree predictors, but the analysis assumed that predictors follow the expected degrees, so is applicable here. It iws conducted by proving, with induction, two relatively simple properties of MPD from which the theorem is proven by induction on the amount of online vertices.

\begin{lemma}
	\label{deg-lm1}
	$p\in[0,1]^n,q\in[0,1]^m$. For some $i\in[n]$, $p'\in[0,1]^{n-1}$ is obtained by removeing entry on index $i$\\
	$\forall k:$\\
	$\Pr[\MPD(I_{p,q})\geq t]\leq\Pr[\MPD(I_{p',q})\geq t-1]$
\end{lemma}
\begin{proof}
	let $V=(v_1,\dots,v_m),U=(u_1,\dots,u_n),U'=U(u_1,\dots,u_{i-1},u_{i+1},\dots,u_n)$\\
	For any bipartite graph $G$ on $(U,V)$ we will generate $G'$ on $(U',V)$ as follows\\
	$\forall v\in V: N_{G'}(v):=N_{G}(v)\cap U'=N_{G}(v)\backslash\set{u_i}$\\
	If $G$ is distributed as $I_{p,q}$ then $G'$ is distributed as $I_{p',q}$. Therefore, all that is necesarry to prove the lemma is to show that $\MPD(G')\geq \MPD(G)-1$, which we will to by induction on $m$\\
	base case: when $m=0$ is trivial. There are no edges so all matchings found on $G,G'$ are of size 0\\
	induction step: will be analyzed in cases.\\
	Let $M_j\subseteq U$ be set of $u\in U$ matched by $\MPD$ on $I$ after $v_1,\dots,v_j$ arrive, $M'_j$ defined symmetrically subset $U'$ on $I'$
	\begin{description}
		\item[Case 1] \MPD doesn't match $u_i$ when running on $I$\\
		In that case the behavior of \MPD on $I$ and $I'$ are identical, $|\MPD(I)|=|\MPD(I')|$, condition holds
		\item[Case 2] \MPD matches $u_i$ to $v_j$ when running on $I$\\
		$\forall j'<j: M_{j'}=M'_{j'}$, particularly for $j'=j-1$
		\begin{description}
			\item[Case 2.1] If $v_j$ is unmatchable on $I'$ then \MPD on both $I,I'$ behave identically from this point onwards, $\MPD(I)=\MPD(I')\cup\set{(u_i,v_j)}$, implying $|\MPD(I)|=|\MPD(I')\cup\set{(u_i,v_j)|=|\MPD(I)|+1}$, condition holds
			\item[Case 2.2] If $v_j$ is matched to $u_{i'}$ on $I'$\\
			Then, $|M_j|=|M'_j|$ and, by the induction hypothesis $|M_m\backslash M_j|\leq|M'_m\backslash M'_j|+1$ {\color{red}(why?)}.\\
			Finally $\MPD(I)=|M_j|+|M_m\backslash M_j|\leq|M'_j|+|M'_m\backslash M'_j|+1=\MPD(I')+1$, condition holds
		\end{description}
	\end{description}
\end{proof}

\begin{lemma}
	\label{deg-lm2}
	$p,p'\in[0,1]^n$ ordered, $q\in[0,1]^m$, $p\leq p'$. \\
	$\forall t:\Pr[\MPD(I_{p,q})\geq t]\leq\Pr[\MPD(I_{p',q})\geq t]$
\end{lemma}
\begin{proof}
	For $S\subseteq [n]$ set of indices of offline vertices, $\Pr_p[N_S]$ is the probability that, in $I_{p,q}$, the first online node's neighborhood is $S$, depending on $p$ and $q_1$\\
	$\Pr_p(N_S)=\prod\limits_{i\in S}p_iq_1\prod\limits_{i\in [n]\backslash S}(1-p_iq_1)$\\
	$p(A,S)$ is vector of ordered weights of the unmatched offline vertices after running algorithm $A$, after the arrival of the first online node\\
	This proof will, once again, rely on indduction on $m$, and the case of $m=0$ is, again, trivial since $\MPD(I_{p,q})=\MPD(I_{p',q})=\emptyset$
	for the induction step:\\
	define $(m-1)$ dimensional vector $q'=(q_2,\dots,q_m)$
	\begin{equation}\begin{aligned}[b]
		\label{deg-eq1}
		\Pr[\MPD(I_{p,q})]\geq t=\Pr_p[N_{\emptyset}]\Pr[\MPD(_{p,q'})\geq t]+\\ \sum\limits_{\emptyset\neq S\subseteq[n]}\Pr_p[N_S]\Pr[\MPD(I_{\vec{p(\MPD,S)},q'})\geq t-1]
	\end{aligned}\end{equation}
	
	let $\vec{r}=pq_1,\vec{r'}=p'q_1$, $\vec{r}_i,\vec{r'}_i$ are probabilities that $v_1$ is connected to $u_i,u'_i$ in $I_{p,q},I_{p',q}$ respectively. $p\leq p'\implies r\leq r'$. Particularly, $1-\vec{r'}_i=(1-s_i)(1-\vec{r}_i)$ for some $s_i$,\\ 
	Let $X_i,Y_i$ independent Bernoulli variables $\Pr[X_i=1]=\vec{r}_i$,$\Pr[Y_i=1]=s_i$.\\
	Let $Z_i=1$ iff $X_i=1\lor Y_i=1$ Bernoulli variable,then $\Pr[Z_i=1]=\vec{r'}_i$\\
	Let $N=\set{i\in[n]|X_i=1}$,$N\supseteq N'=\set{i\in[n]|Z_i=1}$. For any $T\subseteq[n]$.
	$\Delta(S,T):=\Pr[N'=T|N=S]$
	\begin{equation*}
		\Pr_{p'}[N_T]=\Pr[N'=T]=\sum\limits_{S\subseteq T}\Pr[N=S]\Pr[N'=T|N=S]=\sum\limits_{S\subseteq T}\Pr_p(N_S)\Delta(S,T)
	\end{equation*}
	note that $\sum\limits_{T\supseteq S}\Pr[N'=T|N=S]=1$\\
	using \autoref{deg-eq1} we can write
	\begin{equation}\begin{aligned}[b]
		\label{deg-eq2}
		\Pr[\MPD(I_{p',q})\geq t] &= \Pr_{p'}[N_{\emptyset}]\Pr[\MPD(I_{p',q'})\geq t]+\\&
		\sum\limits_{\emptyset\neq T\subseteq[n]}\Pr_{p'}(N_T)\Pr[\MPD(I_{\vec{p'(\MPD,T),q'}})\geq t-1]\\
		&= \Pr_{p'}(N_{\emptyset})\Pr[\MPD(I_{p',q'})\geq t]+ \\&
		\sum\limits_{\emptyset\neq T\subseteq[n]}\sum\limits_{S\subseteq T}\Pr_p[N_s]\Delta(S,T)\Pr[\MPD(I_{\vec{p'(\MPD,T),q'}})\geq t-1]\\
		&=\Pr_p[N_{\emptyset}]\Delta(\emptyset,\emptyset)\Pr[\MPD(I_{p',q'})\geq t]+\\&
		\Pr_p[N_{\emptyset}]\sum\limits_{\emptyset\neq T\subseteq[n]}\Delta(\emptyset,T)\Pr[\MPD(I_{\vec{p'(\MPD,T),q'}})\geq t-1]+\\&
		\sum\limits_{\emptyset\neq S\subseteq[n]}\Pr_p(N_S)\sum\limits_{T\supseteq S}\Delta(S,T)\Pr[\MPD(I_{\vec{p'(\MPD,T),q'}})\geq t-1]
	\end{aligned}\end{equation}
	Note that, for nonempty $S$, $S\subseteq\implies p'(\MPD,S)\leq p'(\MPD,T)$\\
	Using induction hypothesis, for $S,T: T\supseteq S\neq\emptyset$
	\begin{equation}\begin{aligned}[b]
	\Pr[\MPD(I_{\vec{p'(\MPD,T),q'}})\geq t-1]\geq\Pr[\MPD(I_{\vec{p(\MPD,T),q'}})\geq t-1]		
	\end{aligned}\end{equation}
	and
	\begin{equation}
		\label{deg-eq3}
		\Pr[\MPD(I_{p',q'})\geq t]\geq\Pr[\MPD(I_{p,q'})\geq t]
	\end{equation}
	using \autoref{deg-lm1} and induction hypothesis
	\begin{equation}
		\label{deg-eq4}
		\Pr[\MPD(I_{p'(\MPD,T),q'})\geq t-1]\geq\Pr[\MPD(I_{p',q'})\geq t]\geq\Pr[\MPD(I_{p,q'})\geq t]
	\end{equation}
	using \autoref{deg-eq2},\autoref{deg-eq3},\autoref{deg-eq4} with \autoref{deg-eq1} we receive the following
	\begin{align*}
		\Pr[\MPD(I_{p',q})\geq t]\geq \Pr_p[N_\emptyset]\Delta(\emptyset,\emptyset)\Pr[\MPD(I_{p,q'})\geq t]+\\ \Pr_p(N_\emptyset)\sum\limits_{\emptyset\neq T\subseteq[n]}\Delta(\emptyset,T)\Pr(\MPD(I_{p,q'})\geq t)+\\ \sum\limits_{\emptyset\neq S\subseteq[n]}\Pr_p(N_S)\sum\limits_{T\supseteq S}\Delta(S,T)\Pr[\MPD(I_{p(\MPD,S),q'})\geq t-1]
	\end{align*}
	which, when combines with the fact that $\sum\limits_{T\supseteq S}\Delta(S,T)=1$ results in
	\begin{align}
		\Pr[\MPD(I_{p',q})\geq t]\geq\Pr_p[N_\emptyset]\Pr[\MPD(I_{p,q'})\geq t]+\sum\limits_{\emptyset\neq S\subseteq[n]}\Pr_p[N_S]\Pr[\MPD(I_{p(\MPD,S),q'})\geq t-1]=\Pr[\MPD(I_{p,q})\geq t]
	\end{align}
	with the final equality being, again, from \autoref{deg-eq2}
\end{proof}

with these lemmas we can prove \autoref{thm:deg-opt}
\begin{proof}
	We prove, once again, with induction over $m$, and once again the base case $m=0$ is trivial isnce all algorithms return an empty matching with probability of $1$\\
	for the induction step, when inequality holds for $m-1$ online nodes:\\
	$q'=(q_2,\dots,q_m)$, $2^{[n]}\supseteq R_A=\set{S\subseteq[n]|p(A,S)=p}$
	by induction hypothesis
	\begin{align*}
		\Pr[A(I_{p,q})\geq t]&=\sum\limits_{S\in R_A}\Pr_p(N_S)\Pr[A(I_{p,q'})\geq t]+\\
		&\sum\limits_{S\in[n]\backslash R_A}\Pr_p[N_S]\Pr[A(I_{p(A,S),q'})\geq t-1]\\
		&\leq\sum\limits_{S\in R_A}\Pr_p(N_S)\Pr[\MPD(I_{p,q'})\geq t]+\\
		&\sum\limits_{S\in[n]\backslash R_A}\Pr_p[N_S]\Pr[\MPD(I_{p(A,S),q'})\geq t-1]
	\end{align*}
	$S\in R_A\implies p(A,S)\leq p(A_0,S)$. Applying \autoref{deg-lm1} to the first term, \autoref{deg-lm2} to the second term results in
	\begin{align*}
		&\Pr[A(I_{p,q})]\leq\Pr_p[N_\emptyset]\Pr[\MPD(I_{p,q'})\geq t]+\\
		&\sum\limits_{\emptyset\neq S\in R_A}\Pr_p[N_S]\Pr[\MPD(I_{p(\MPD,S),q'})\geq t-1]+\\
		&\sum\limits_{S\subseteq[n]\backslash R_A}\Pr_p[N_S]\Pr[\MPD(I_{p(\MPD,S),q'})\geq t-1]\\
		=&\Pr_p[N_\emptyset]\Pr[\MPD(I_{p,q'})\geq t]+\\
		\sum\limits_{\emptyset\neq S\subseteq[n]}\Pr_p[N_S]\Pr[\MPD(I_{p(\MPD,S),q'})\geq t-1]\\
		=&\Pr(\MPD(I_{p,q'})\geq t)
	\end{align*}
\end{proof}